\DocumentMetadata{
  pdfversion=2.0,
  pdfstandard=ua-2,
  lang=en-US,
  tagging=on,
  tagging-setup={math/setup=mathml-SE,tabsorder=structure}
}
\documentclass[11pt,letterpaper]{article}
\usepackage[margin=0.68in,includefoot]{geometry}
\usepackage{newcomputermodern}
\usepackage{amsmath,amscd,mathtools}
\usepackage{tikz,adjustbox}
\providecommand{\square}{\mdlgwhtsquare}
\usepackage[hidelinks]{hyperref}
\hypersetup{
  pdftitle={Numerical Analysis PhD Exam, January 2024},
  pdfauthor={University of Florida Department of Mathematics},
  pdfsubject={Graduate mathematics examination},
  pdfdisplaydoctitle=true,
  bookmarksopen=true
}
\setcounter{secnumdepth}{0}
\setlength{\parindent}{0pt}
\setlength{\parskip}{0.28em}
\setlength{\emergencystretch}{3em}
\setlength{\leftmargini}{2.15em}
\setlength{\leftmarginii}{2.65em}
\setlength{\leftmarginiii}{3em}
\renewcommand{\labelenumi}{\textbf{\theenumi.}}
\pagestyle{empty}
\raggedbottom
\allowdisplaybreaks
\newenvironment{examproblems}[1]{%
  \begin{enumerate}%
  \setcounter{enumi}{#1}%
  \setlength{\topsep}{0.35em}%
  \setlength{\partopsep}{0pt}%
  \setlength{\itemsep}{0.48em}%
  \setlength{\parsep}{0.18em}%
}{\end{enumerate}}
\newenvironment{examsubparts}{%
  \begin{enumerate}%
  \setlength{\topsep}{0.18em}%
  \setlength{\partopsep}{0pt}%
  \setlength{\itemsep}{0.16em}%
  \setlength{\parsep}{0.1em}%
}{\end{enumerate}}
\newenvironment{examinstructions}{%
  \par\begingroup\itshape
}{%
  \par\endgroup\vspace{0.18em}
}
\makeatletter
\renewcommand\section{\@startsection{section}{1}{0pt}%
  {-1.25ex plus -.25ex minus -.1ex}%
  {0.55ex plus .1ex}%
  {\normalfont\large\bfseries}}
\renewcommand{\@maketitle}{%
  \newpage\null\vskip -1em%
  {\footnotesize\bfseries
    \href{https://www.ufl.edu/}{University of Florida}, \href{https://math.ufl.edu/}{Department of Mathematics}\par}%
  \vskip -0.5em%
  \tagstructbegin{tag=Title}%
    {\LARGE\bfseries\href{https://gma.math.ufl.edu/exams/numerical-analysis-phd/}{Numerical Analysis PhD Exam}, January 2024\par}%
  \tagstructend%
  \vskip 0.25em%
  \tagmcbegin{artifact}%
    \hrule height 1pt%
  \tagmcend%
  \vskip 1em%
}
\makeatother
\title{Numerical Analysis PhD Exam, January 2024}
\author{}
\date{}
\begin{document}
\maketitle
\thispagestyle{empty}
\begin{examinstructions}
Do 4 (four) of the first 5 (1–5) and 4 (four) of the last 5 problems (6–10).
\end{examinstructions}
\begin{examproblems}{0}
\item
Assume \(A \in \mathbb{C}^{m \times m}\).
\begin{examsubparts}
\item[\textnormal{(a)}]
Show that~\(A\) has a Schur decomposition.
\item[\textnormal{(b)}]
If~\(A\) has a collection of~\(m\) linearly independent eigenvectors, show that~\(A\) is diagonalizable.
\end{examsubparts}
\item
Let matrix \(A \in \mathbb{C}^{m \times n}\) with \(n < m\). Let \(b \in \mathbb{C}^m\), and let~\(r\) denote the residual vector \(r=b-Ax\).
\begin{examsubparts}
\item[\textnormal{(a)}]
Show that~\(x\) solves the least squares problem \(\min \lVert b-Ax\rVert_2\) if and only if \(r \in \operatorname{Null}(A^*)\).
\item[\textnormal{(b)}]
Suppose~\(A\) is full rank, and describe how to find the least squares solution using the QR decomposition of~\(A\).
\end{examsubparts}
\item
Let \(A \in \mathbb{C}^{m \times n}\), with \(m \geq n\) and \(\operatorname{rank}(A)=p=n \geq 3\). Let \(a_1,a_2,\cdots\) denote the columns of~\(A\).
\begin{examsubparts}
\item[\textnormal{(a)}]
Using the modified Gram–Schmidt process, write out expressions for \(q_1,q_2,q_3\), the first three columns of~\(Q\) in the QR decomposition of~\(A\).
\item[\textnormal{(b)}]
Show the vector \(q_3\) found in part (a) is orthogonal to both \(q_1\) and \(q_2\).
\end{examsubparts}
\item
Let \(\lVert\,\cdot\,\rVert\) be a subordinate (induced) matrix norm.
\begin{examsubparts}
\item[\textnormal{(a)}]
If~\(E\) is a \(n \times n\) with \(\lVert E\rVert<1\), then show \(I+E\) is nonsingular and 
\[
\lVert(I+E)^{-1}\rVert \leq \frac{1}{1-\lVert E\rVert}.
\]
\item[\textnormal{(b)}]
If~\(A\) is a \(n \times n\) invertible and~\(E\) is \(n \times n\) with \(\lVert A^{-1}\rVert\lVert E\rVert<1\), then show \(A+E\) is nonsingular and 
\[
\lVert(A+E)^{-1}\rVert \leq \frac{\lVert A^{-1}\rVert}{1-\lVert A^{-1}\rVert\lVert E\rVert}.
\]
\end{examsubparts}
\item
If \(q_1,\ldots,q_n\) is an orthonormal basis for the subspace \(V \subset \mathbb{C}^m\) with \(m>n\), prove that the orthogonal projector onto~\(V\) is \(QQ^*\), where~\(Q\) is the matrix whose columns are the \(q_j\).
\item
Derive the three-point formula for the second derivative 
\[
f''(x_0)=\frac{1}{h^2}\bigl(f(x_0-h)-2f(x_0)+f(x_0+h)\bigr)-\frac{h^2}{12}f^{(4)}(\eta),
\]
 for some \(\eta \in [x_0-h,x_0+h]\).
\item
Consider 
\[
y'=4ty; \qquad t\in[0,1]; \qquad y(0)=2, \tag{1}
\]
 which has solution \(Y(t)=2e^{2t^2}\).
\begin{examsubparts}
\item[\textnormal{(a)}]
Derive an error bound for the forward Euler scheme.
\item[\textnormal{(b)}]
Derive the Taylor method of order 2 for (1).
\end{examsubparts}
\item
Let \(P_1\) be the space of polynomials of degree at most one. Using the norm 
\[
\lVert u\rVert_2=\left(\int_a^b u^2\,dx\right)^{1/2}.
\]
\begin{examsubparts}
\item[\textnormal{(a)}]
Find the least-squares approximation to \(f(x)=x^3\) in \(P_1\) over \([a,b]=[-1,1]\).
\item[\textnormal{(b)}]
Find the least-squares approximation to \(f(x)=x^4\) in \(P_1\) over \([a,b]=[0,1]\).
\end{examsubparts}
\item
Consider the fixed point problem \(x=f(x)\) where \(f(x)=e^{-(3+x)}\).
\begin{examsubparts}
\item[\textnormal{(a)}]
Assuming all computations are done in exact arithmetic, find the largest open interval in~\(\mathbb{R}\) where the fixed point iteration \(x_{k+1}=f(x_k)\) is ensured to converge.
\item[\textnormal{(b)}]
Write a Newton iteration for finding the fixed point.
\end{examsubparts}
\item
Suppose \(f \in C^{n+1}[a,b]\), and let \(p \in P_n\) be a polynomial that interpolates the data \(\{(x_i,f(x_i))\}_{i=0}^n\), where \(x_0,\ldots,x_n\) are distinct points in \([a,b]\). Consider an arbitrary fixed \(x \in [a,b]\), and derive an exact expression for the error \(f(x)-p(x)\).
\end{examproblems}
\end{document}
